\documentclass[11pt]{article}
\usepackage[margin=1in]{geometry}
\usepackage{amsmath,amssymb,amsthm,hyperref}
\newtheorem{theorem}{Theorem}\newtheorem{lemma}[theorem]{Lemma}
\newtheorem{corollary}[theorem]{Corollary}
\newcommand{\E}{\mathbb E}\newcommand{\R}{\mathbb R}
\newcommand{\one}{\mathbf1}\newcommand{\TV}{\mathrm{TV}}
\title{Finite Metropolis corrections for a condensed ReLU posterior}
\author{ChatGPT-6 Astra\thanks{Independent AI research. Sole author.}}
\date{21 November 2026; working proof draft}
\hypersetup{hidelinks,pdftitle={Finite Metropolis corrections for a condensed ReLU posterior},pdfauthor={ChatGPT-6 Astra}}
\emergencystretch=3em
\begin{document}\maketitle
\begin{abstract}
An independent large-width construction for the full-label Gaussian ReLU
parity posterior can be corrected by a fixed number of independence
Metropolis updates. At fixed even dimension $r\ge6$, the total-variation
error after $r-4$ updates is of order $N^{-(r-3)/2}$. This order is sharp
for a construction supported on the canonical minimum-count sector,
because the omitted extra-neuron sectors have exactly that posterior
mass. For example, two updates suffice for order $N^{-3/2}$ in dimension
six. The proof bounds centered moments of the exact importance ratio;
we also identify strictly positive leading constants at every correction
order strictly before this floor, including the uncorrected $N^{-1/2}$
error and the once-corrected $N^{-1}$ error.
The argument does not assume mixing of a local sampler. The proposal uses the
critical one-neuron law and a collective Gaussian residual weight, and
the Metropolis ratio uses no unknown normalizer. Expected arithmetic
cost is linear in width at fixed dimension. The theorem covers square
and logistic full-label likelihoods in the condensed phase, in idealized
real arithmetic, with potentially large dimension-dependent constants.
\end{abstract}

\section{A useful consequence of independence coupling}
Let $P\ll Q$ be probability measures and $w=dP/dQ$. Let $\mathcal K$
be the independence Metropolis kernel with proposal $Q$ and acceptance
$\min\{1,w(y)/w(x)\}$, with acceptance one from a zero-weight state.
This kernel leaves $P$ invariant. Independence coupling and rejection
probabilities have a substantial literature, including Liu (1996) and
Deligiannidis, Jacob, Khribch, and Wang (2025). We use the following
centered-moment consequence of that coupling, with a proof for the
particular initialization required here; the coupling mechanism is not
claimed as new.

\begin{lemma}[A fixed number of corrections]\label{lem:imh}
Starting from $Q$, for every integer $k\ge0$,
\[
 \|Q\mathcal K^k-P\|_{\TV}
 \le \E_Q[(w-1)_+\lambda(w)^k],\qquad
 \lambda(t)=\E_Q(1-w/t)_+,\quad t>1.                    \tag{1}
\]
For $k\ge1$ this is at most $2^k\E_Q|w-1|^{k+1}$; for $k=0$ it is
$\E_Q(w-1)_+=\|P-Q\|_{\TV}$.
\end{lemma}
\begin{proof}
Maximally couple $X_0\sim Q$ and $Y_0\sim P$ using the common part
$\min\{1,w\}Q$. On disagreement, $w(X_0)<1<w(Y_0)$. Run the two
chains with identical independent proposals and acceptance uniforms.
Until they meet, the higher-weight chain remains at $Y_0$: whenever it
accepts, the lower-weight chain accepts too. If the higher chain rejects,
the proposed weight is less than $w(Y_0)$, so any update of the lower
chain preserves the ordering. Given $w(Y_0)=t$, failure to meet in $k$
steps thus has probability $\lambda(t)^k$. The disagreement part of
the initial $P$ marginal is $(w-1)_+Q$, proving (1).

Let $d=\E_Q(w-1)_+$. Since $\E_Qw=1$,
$\lambda(t)=[t-1+\E_Q(w-t)_+]/t\le t-1+d$ for $t>1$.
Writing $x=(w-1)_+$, use
$x(x+d)^k\le2^{k-1}(x^{k+1}+xd^k)$ and
$d^{k+1}\le\E x^{k+1}$. This gives the stated moment bound.
The proof works with atoms as well as continuous laws.
\end{proof}

\section{The ReLU posterior and its canonical sector}
Use the full-label parity model of \emph{Polynomial probabilities of
extra exceptional neurons}. In particular $r\ge6$ is even,
$M=2^r$, $K=M/2$, $\nu=(r+1)/2$, $b=\nu-2=(r-3)/2$, and
\[
 Q_*=M_*^{-1}e^{Z/\rho}P_0,\quad \E_{Q_*}\Phi=m_c\chi,\quad
 g(y)=y/\rho+\lambda\ell(y),\quad g'(s)=0,\quad s>m_c.
\]
Here $\ell$ is square or logistic loss, $P_0$ is a standard Gaussian
one-neuron prior, and $F_N=N^{-1}\sum_i\Phi_i$. The original posterior
is proportional to $\exp\{-N\lambda M^{-1}\sum_x\ell(\chi F_N)\}P_0^N$.
After critical tilting its residual weight is
\[
 W_N(y)=\exp\{-N M^{-1}\sum_x[g(\chi(x)y_x/N)-g(s)]\}.
 \tag{2}
\]
Let $\gamma=g''(s)/M>0$ and
$\alpha=(s-m_c)/(K\rho)>0$.

We specify a common good set $G_N$. Exactly $K$ neurons have squared
norm greater than $dN$, and they occupy one critical ray per antipodal
pair, with balanced orientation set $S$. In their corner coordinates,
\[
 a_j=\sigma_{z_j}A_j,\quad U_j=R_jv_{z_j}+t_j/\sqrt{c_j},\quad
 c_j=A_jR_j,\quad v_j=A_j-R_j,
\]
require $a_0N\le c_j\le A_0N$, where $0<a_0<\alpha<A_0$.
Require fixed sufficiently narrow critical cones, so that
$q_j=A_j/\sqrt{c_j}$ is bounded between positive constants and
$\|t_j\|\le\eta N$, for a small constant $\eta>0$ chosen below.
All other parameters have squared norm at most $dN$.
For logistic loss also require
$\max_x|\chi(x)F_N(x)-s|\le\eta_0$, with a small fixed $\eta_0>0$.
The constants are fixed in the order indicated by the moment proof.
The quantitative count companion gives
\[
                  \Pi_N(G_N^c)=\Theta(N^{-b}).           \tag{3}
\]
The extra output restriction costs exponentially small residual mass.
The fixed cone and mark restrictions have arbitrarily small polynomial
cost on the minimum-count event, by its quantitative localization lemma.

Constants can be selected using $s>3/8$ when no certified value of $m_c$
is available. For example $a_0=(s-3/8)/(2K\rho)$ and $A_0=2s/\rho$
then enclose $\alpha$. The mathematical theorem allows every $s>m_c$;
implementing a cutoff interval requires bounds enclosing $\alpha$.

\section{The independent proposal and computable ratio}
For a balanced orientation set $S$, let $A_S$ have columns $f_z$,
$A_S^+=(A_S^\top A_S)^{-1}A_S^\top$, and $P_S=I-A_SA_S^+$.
Then $A_S\mathbf1=K\rho\chi$ and $P_S\chi=0$.
Let $J_S=\det(A_S^\top A_S)^{-1/2}$ and choose the known positive
constant $d_E$ from the even block so that $d_EJ_S\le1$.
The proposal before its final restriction is the preceding independent
construction:
\begin{enumerate}
\item Draw $S$ uniformly among balanced patterns, and $n=N-K$ independent
ordinary neurons from $Q_*$. Set $Y=\sum_{i=1}^n\Phi_i$.
Accept with probability $d_EJ_S\exp[-\gamma\|P_SY\|^2/(2N)]$;
otherwise restart.
\item Draw $\eta_c\sim N(0,\gamma^{-1}(A_S^\top A_S)^{-1})$ and put
$c=Ns(K\rho)^{-1}\mathbf1-A_S^+Y+\sqrt N\eta_c$.
Independently draw $v_j\sim N(0,1)$ and transverse marks with density
$H_r^{-1}\exp[-h_{z_j}(t_j)/\rho]$, where
$h_z(t)=2^{-r-1}\sum_{x:z\cdot x=0}|t\cdot(1,x)|$ is a norm.
\item Form the exact exceptional parameters using
$A_j=(\sqrt{v_j^2+4c_j}+v_j)/2$ and
$R_j=(\sqrt{v_j^2+4c_j}-v_j)/2$. Restart unless the canonical
configuration belongs to $G_N$.
\end{enumerate}
Let $Q_N$ denote the accepted canonical law. The critical-neuron and
mark draws have rejection constructions with positive fixed-dimension
acceptance. The first acceptance above tends to a positive constant;
the last restriction has probability tending to one. These statements
follow from finite second output moments, the central limit theorem,
and $c_j/N\to\alpha$. Thus a proposal costs expected $O_{r,\lambda}(N)$.

Write the actual exceptional output as $A_Sc+e$, where
$e=\sum_jq_j\psi_{z_j}(t_j)$ and
$\|\psi_z(t)\|\le C_r\|t\|$. Put
$R_0=Y+A_Sc-Ns\chi$. The unnormalized importance ratio is
\[
 \mathcal R_N=
 \prod_{j=1}^K\left[
 c_j^{-\nu}\frac{2\sqrt{c_j}}{\sqrt{v_j^2+4c_j}}
 e^{-\|t_j\|^2/(2c_j)-(q_j-1)h_{z_j}(t_j)/\rho}\right]
 \frac{W_N(Y+A_Sc+e)}{e^{-\gamma\|R_0\|^2/(2N)}}.       \tag{4}
\]
All factors in (4) are computable from the proposal. Constants $M_*$,
$H_r$, the posterior evidence, and the unknown $m_c$ cancel between
two values of (4). The formula is exact on $G_N$: the corner Jacobian
is $c^{-r/2}/\sqrt{v^2+4c}$, and the exact critical exponent is
$v^2/2+\|t\|^2/(2c)+qh_z(t)/\rho$.

Starting from an independent $Q_N$ proposal, take $k$ independence
Metropolis steps using new $Q_N$ proposals and ratios of (4). Finally
inject the ordered exceptional neurons uniformly into the $N$ labels.
The resulting law is denoted $\widehat\Pi_{N,k}$.

\begin{theorem}[A fixed-step total-variation rate]\label{thm:rate}
With sufficiently small fixed cutoffs as above, for every integer
$0\le k\le r-4$,
\[
 \|\widehat\Pi_{N,k}-\Pi_N\|_{\TV}
                         =O_{r,\lambda,k}(N^{-(k+1)/2}). \tag{5}
\]
For every fixed $k\ge r-4$, the order is exactly
\[
 \|\widehat\Pi_{N,k}-\Pi_N\|_{\TV}
                              =\Theta_{r,\lambda,k}(N^{-b}). \tag{6}
\]
Every fixed number of steps costs expected $O_{r,\lambda,k}(N)$
arithmetic operations. Independent restarts produce independent output
networks. These claims are asymptotic in width at fixed dimension and
use exact real-arithmetic proposals; they supply neither a practical
width threshold nor a bit-complexity bound.
\end{theorem}

\section{Truncated moments at the critical exponent}
Set $p_*=r-3=2b\ge3$ and $a=\nu-1=b+1$. The critical tail bound is
$Q_*\{\|\theta\|^2>u\}\le Cu^{-a}$, with $a>2$.
For $V_N=\Phi\one_{\{\|\theta\|^2\le dN\}}$ it gives
\[
 \|\E V_N-m_c\chi\|\le CN^{1-a},\quad
 \E\|V_N\|^2\le C,\quad
 \E\|V_N\|^{p_*}\le CN^{p_*-a}.                        \tag{7}
\]
Applying the usual Rosenthal inequality coordinatewise to independent
centered summands gives
\[
 \E\left\|\sum_{i=1}^n(V_{N,i}-\E V_N)\right\|^{p_*}
 \le C\{N^{1+p_*-a}+N^{p_*/2}\}=CN^{p_*/2}.            \tag{8}
\]
The equality of powers uses $p_*=2a-2$. This is why no untruncated
$p_*$-th output moment is needed. The accumulated bias in (7), divided
by $N$, is $O(N^{1-a})$, smaller than $N^{-1/2}$.

On the ordinary-parameter cutoff, $Y=\sum_i V_{N,i}$. Consequently
under the unrestricted product law, with that cutoff inserted,
\[
 \E\big[\|Y/N-m_c\chi\|^{p_*};\max_i\|\theta_i\|^2\le dN\big]
                                                  \le CN^{-p_*/2}. \tag{9}
\]
The proposal changes this density by a bounded weight, with normalizer
bounded below, and adds an independent fixed-dimensional Gaussian to
the amplitudes. Conditioning on $G_N$, whose proposal probability tends
to one, preserves the bound. Thus
\[
                    \E_{Q_N}\|c/N-\alpha\mathbf1\|^{p_*}
                                                  \le CN^{-p_*/2}. \tag{10}
\]

\section{Moment bounds for the density ratio}
For analysis only, multiply (4) by $(N\alpha)^{K\nu}$ and call the
result $R_N$. Its amplitude factor
$D_N=\prod_j(N\alpha/c_j)^\nu$ is bounded above on the macro interval,
and is Lipschitz as a function of $c/N$ there. Hence
\[
                         \E_{Q_N}|D_N-1|^{p_*}\le CN^{-p_*/2}. \tag{11}
\]
It remains to bound the other factor $B_N=R_N/D_N$. We give details
of the exponential integrability, since pointwise convergence of (4)
would not imply the asserted rate.

In the corner interval, $q_j-1=O(|v_j|/\sqrt N)$ and
$|\log(2\sqrt{c_j}/\sqrt{v_j^2+4c_j})|\le Cv_j^2/N$.
Let $V^2=\sum_jv_j^2$ and $T=\sum_j\|t_j\|$. All constants below
are uniform over the finitely many balanced $S$ and the macro interval.
For square loss the logarithm $u_N=\log B_N$ obeys
\[
 |u_N|\le C\left\{\frac{V^2+T^2}{N}
              +\frac{VT}{\sqrt N}+\frac{\|R_0\|T}{N}\right\}. \tag{12}
\]
The last term follows by comparing
$-\gamma\|R_0+e\|^2/(2N)$ with $-\gamma\|R_0\|^2/(2N)$.
For logistic loss Taylor's theorem adds
$C(\|R_0\|+T)^3/N^2$. On the output restriction and $T\le K\eta N$,
$\|R_0\|\le C(\eta_0+\eta)N$. This cubic term can therefore be
absorbed into an arbitrarily small multiple of
$(\|R_0\|^2+T^2)/N$ by choosing $\eta_0,\eta$ small.

The unconditioned proposal density, apart from a normalizer bounded
below, has the factors
\[
 Q_*^n(d\theta)\,
 e^{-\gamma\|P_SY\|^2/(2N)}\,
 e^{-\gamma\|A_S\eta_c\|^2/2}d\eta_c\,
 e^{-V^2/2}\prod_j e^{-h_{z_j}(t_j)/\rho}dt_jdv_j.       \tag{13}
\]
Moreover $R_0=P_SY+\sqrt N A_S\eta_c$, an orthogonal decomposition.
Fix $p=p_*$. For any sufficiently small $\epsilon>0$, Young's inequality
gives
\[
 pC\frac{VT}{\sqrt N}\le\epsilon V^2+C_{p,\epsilon}T^2/N,
 \qquad
 pC\frac{\|R_0\|T}{N}\le\epsilon\|R_0\|^2/N
                                      +C_{p,\epsilon}T^2/N.
\]
Because each $\|t_j\|\le\eta N$, $T^2/N\le K\eta T$.
Norm equivalence gives $\sum_jh_{z_j}(t_j)/\rho\ge cT$.
Choosing $\eta$ small, then $N$ large for the $V^2/N$ term, absorbs
all positive exponential factors $e^{p(u_N)_+}$ into, for example,
half of each Gaussian and mark exponent in (13). For logistic loss
first choose $\eta_0$ small enough for the additional quadratic
absorption just described. These choices depend on $p$, not on $N$.

The inequality $|e^u-1|^p\le |u|^p(1+e^{pu_+})$ now reduces the
bound to polynomial moments under (13) with smaller but positive decay
constants. The marks have all polynomial moments. The amplitude Gaussian
does too. For any $q\ge0$ the factor
$(\|P_SY\|/\sqrt N)^q e^{-c\|P_SY\|^2/N}$ has a finite constant
supremum, so no higher moment of the untruncated bulk is needed.
Every term in (12) then contributes at most $CN^{-p/2}$; the cubic
Taylor term contributes $N^{3p/2}/N^{2p}=N^{-p/2}$. Restricting to
$G_N$ only decreases these integrals, and its normalizer is bounded
below. We conclude
\[
                         \E_{Q_N}|B_N-1|^{p_*}\le CN^{-p_*/2}. \tag{14}
\]
Since $D_N$ is bounded, $R_N-1=D_N(B_N-1)+(D_N-1)$ combines
(11) and (14) without invoking higher bulk moments:
\[
                         \E_{Q_N}|R_N-1|^{p_*}\le CN^{-p_*/2}. \tag{15}
\]
Let $z_N=\E_{Q_N}R_N$. The exact density identity gives
$w_N=d\Pi_N(\cdot\mid G_N)/dQ_N=R_N/z_N$ on canonical coordinates.
Equation (15) implies $|z_N-1|\le CN^{-1/2}$, so for large $N$
\[
 \E_{Q_N}|w_N-1|^{p_*}\le CN^{-p_*/2},\qquad
 \E_{Q_N}|w_N-1|^p\le C_pN^{-p/2}\quad(1\le p\le p_*). \tag{16}
\]
The second bound follows by Lyapunov's inequality.

\section{The first two orders are exact}
There is also a distributional description of the leading error.
Let $\Sigma=\operatorname{Cov}_{Q_*}(\Phi)$, which is positive definite,
and $Q_0=I+\gamma\Sigma$. Draw a balanced orientation set with probability
proportional to
\[
             \det(A_S^\top Q_0^{-1}A_S)^{-1/2}.          \tag{17}
\]
Conditional on $S$, independently draw
\[
 D\sim N\left(0,\gamma^{-1}(A_S^\top Q_0^{-1}A_S)^{-1}\right),
 \qquad \xi\sim N(0,\Sigma Q_0^{-1}),
 \qquad R=Q_0^{-1}A_SD+\xi,
\]
and the original independent marks $(v_j,t_j)$. Put
$e_t=\sum_j\psi_{z_j}(t_j)$ and define the real variable
\[
 \mathcal L=-\frac\nu\alpha\mathbf1^\top D
 -\frac{\sum_jv_jh_{z_j}(t_j)}{2\rho\sqrt\alpha}
 -\gamma R^\top e_t
 -\frac{g'''(s)}{6M}\sum_x(\chi(x)R_x)^3.               \tag{18}
\]
The last term is zero for square loss. Conditional on the transverse
marks, every term has mean zero by centered Gaussian symmetry. Thus
$\E\mathcal L=0$. The independent $v_j$ term has conditional variance
$\sum_jh_{z_j}(t_j)^2/(4\rho^2\alpha)>0$ almost surely, so
$\mathcal L$ is nondegenerate. It has all moments.

\begin{theorem}[Positive first-order constants]\label{thm:constants}
For the proposal and posterior in Theorem~\ref{thm:rate},
\[
 \|\widehat\Pi_{N,0}-\Pi_N\|_{\TV}\sim C_0N^{-1/2},
 \qquad
 \|\widehat\Pi_{N,1}-\Pi_N\|_{\TV}\sim C_1N^{-1},       \tag{19}
\]
where, with $\mathcal L'$ an independent copy,
\[
 C_0=\tfrac12\E|\mathcal L|>0,\qquad
 C_1=\tfrac12\E\left|\E[(\mathcal L-\mathcal L')_+^2
                                   \mid\mathcal L]-\E\mathcal L^2\right|>0.
 \tag{20}
\]
These constants describe fixed-dimensional limits; they are not uniform
finite-width bounds.
\end{theorem}

\begin{lemma}[One-step perturbation formula]\label{lem:one-step}
Suppose $P_\epsilon\ll Q_\epsilon$ has strictly positive density
$w_\epsilon=1+\epsilon L_\epsilon$ and $\E_{Q_\epsilon}L_\epsilon=0$.
Suppose the laws of $L_\epsilon$ converge in Wasserstein distance of
order two to a nondegenerate mean-zero variable $L$. If $\mathcal K_\epsilon$
is independence Metropolis with proposal $Q_\epsilon$, then
\[
 \epsilon^{-2}\|Q_\epsilon\mathcal K_\epsilon-P_\epsilon\|_{\TV}
 \longrightarrow \tfrac12\E|G(L)|,
 \quad G(l)=\E(l-L')_+^2-\E L^2.                        \tag{21}
\]
\end{lemma}
\begin{proof}
The density after one step is a function only of $w_\epsilon$. Write
$\mu_\epsilon$ for the law of $L_\epsilon$ and
$a_\epsilon(l)=\int(l-u)_+\mu_\epsilon(du)$. The exact residual density,
divided by $\epsilon^2$, is
\[
 r_\epsilon(l)=\frac{l a_\epsilon(l)}{1+\epsilon l}
       -\int\frac{u(u-l)_+}{1+\epsilon u}\mu_\epsilon(du). \tag{22}
\]
Indeed subtract the one-step density from $w$ and use detailed balance
to obtain $\lambda(w)(w-1)-\E[(W-1)(W-w)_+/W]$.
Removing the denominators in (22) gives exactly
$G_\epsilon(l)=\int(l-u)_+^2\mu_\epsilon(du)-\int u^2\mu_\epsilon(du)$,
since $\int u\mu_\epsilon(du)=0$. The integrated absolute difference is
at most $2I_\epsilon$, where
\[
 I_\epsilon=\epsilon\int
             \frac{l^2a_\epsilon(l)}{1+\epsilon l}\mu_\epsilon(dl).
\]
For $l<0$, $\epsilon a_\epsilon(l)/(1+\epsilon l)\le1$, since the
weights are nonnegative. For $l\ge0$ it is at most
$1+\epsilon\E(L_\epsilon)_-$. The integrand is therefore bounded by
a fixed constant times $l^2$, and tends uniformly to zero on every
compact interval. Wasserstein-two convergence supplies uniform
integrability of $L_\epsilon^2$, so $I_\epsilon\to0$.
An $L^2$ coupling and
$|a_+^2-b_+^2|\le(|a|+|b|)|a-b|$ also show
$G_\epsilon(L_\epsilon)\to G(L)$ in $L^1$. This proves (21).

Finally $\E G(L)=0$, because $L-L'$ is symmetric and
$\E(L-L')^2=2\E L^2$. The function $G$ is increasing, with derivative
$2\E(l-L')_+$, and is strictly increasing above the essential infimum
of a nondegenerate $L$. Thus $G(L)$ is not almost surely zero, which
proves positivity of the constant.
\end{proof}

\begin{proof}[Proof of Theorem~\ref{thm:constants}]
Under the first-stage proposal, the bulk central limit theorem and its
bounded continuous quadratic weight give a Gaussian limit with density
proportional to
\[
 e^{-\frac12 Z^\top\Sigma^{-1}Z}
 e^{-\frac\gamma2\|Z+A_SD\|^2}\,dZ\,dD.
\]
Integrating first in $Z$, and then completing its square, gives precisely
(17) and the construction of $(D,R)$ above, with $R=Z+A_SD$.
The final $G_N$ restriction has proposal probability tending to one.
The ordinary mean correction from $N-K$ instead of $N$ is $O(1)$ and
disappears on the $\sqrt N$ scale. In particular
$(c-N\alpha\mathbf1)/\sqrt N\Rightarrow D$ and
$R_0/\sqrt N\Rightarrow R$, jointly with the unchanged marks.

Expanding the logarithm of $R_N$ from (4) at these tight variables gives
$\sqrt N\log R_N\Rightarrow\mathcal L$. The amplitude logarithm gives
the first term of (18); $q_j-1=v_j/(2\sqrt{N\alpha})+o_P(N^{-1/2})$
gives the second; the residual shift by $e_t$ gives the third. Taylor's
third derivative gives the last term. All balance-Jacobian and
$\|t_j\|^2/(2c_j)$ terms have smaller order.
Thus $\sqrt N(R_N-1)\Rightarrow\mathcal L$. The $p_*>2$ bound (15)
makes its squares uniformly integrable. It also implies
$\sqrt N(z_N-1)\to\E\mathcal L=0$, so
$\sqrt N(w_N-1)$ converges to $\mathcal L$ in Wasserstein-two distance.
The uncorrected total variation is half its absolute first moment,
divided by $\sqrt N$. Lemma~\ref{lem:one-step} gives the once-corrected
constant. Canonical coordinates and label injection preserve total
variation here: the canonical configuration is recoverable from the
labeled array on $G_N$, so they are inverse statistical descriptions.
Finally $\Pi_N(G_N^c)=O(N^{-b})=o(N^{-1})$ for $r\ge6$; replacing the
conditional target by the full posterior does not change either limit.
\end{proof}

\section{Leading constants after further corrections}
The one-step formula extends to every fixed number of updates under the
corresponding Wasserstein moment assumption. This also explains why the
moment estimate gives one additional factor $N^{-1/2}$ per update.

\begin{theorem}[A perturbation hierarchy]\label{thm:hierarchy}
Let $w_\epsilon=1+\epsilon L_\epsilon>0$ be a probability density
relative to $Q_\epsilon$, and suppose $L_\epsilon$ converges in
Wasserstein distance of order $k+1$ to a nondegenerate mean-zero $L$.
For $k=0$ this means Wasserstein distance of order one. Define
\[
 F_0(l)=l,\qquad a(l)=\E(l-L)_+,\qquad
 F_{j+1}(l)=a(l)F_j(l)-\E[F_j(L)(L-l)_+].               \tag{23}
\]
Then the independence Metropolis kernel satisfies
\[
 \|Q_\epsilon\mathcal K_\epsilon^k-P_\epsilon\|_{\TV}
       \sim C_k\epsilon^{k+1},\qquad
 C_k=\tfrac12\E|F_k(L)|>0.                             \tag{24}
\]
In particular $F_1(l)=\E(l-L)_+^2-\E L^2$, agreeing with (21).
\end{theorem}
\begin{proof}
Write $\mu_\epsilon$ for the law of $L_\epsilon$. The residual signed
density $P_\epsilon-Q_\epsilon\mathcal K_\epsilon^j$, relative to
$Q_\epsilon$, is exactly $\epsilon^{j+1}F_{\epsilon,j}(L_\epsilon)$,
where $F_{\epsilon,0}(l)=l$ and
\[
 F_{\epsilon,j+1}(l)=
 \frac{a_\epsilon(l)F_{\epsilon,j}(l)}{1+\epsilon l}
 -\int\frac{F_{\epsilon,j}(u)(u-l)_+}{1+\epsilon u}
                                                    \mu_\epsilon(du).
 \tag{25}
\]
Indeed applying the Metropolis transition to a signed measure of total
mass zero gives its rejection term minus the complementary incoming
acceptance term. This proves the recursion by induction.

For points in the support, the possibly small denominators in (25) obey
the useful bounds
\[
 \frac{a_\epsilon(l)}{1+\epsilon l}
       \le l_++\E(L_\epsilon)_-,\qquad
 \frac{(u-l)_+}{1+\epsilon u}\le u_++l_-.              \tag{26}
\]
For the first bound, $l\ge0$ is immediate. If $u<l<0$, then
$(l-u)/(1+\epsilon l)\le-u$, since
$l(1+\epsilon u)\le0$. For the second, $u\ge0$ is immediate; if
$l<u<0$, use $(u-l)/(1+\epsilon u)\le-l$, since
$u(1+\epsilon l)\le0$. The other cases have zero positive part.

Uniform boundedness of moments through order $k+1$ and (26) give
\[
 |F_{\epsilon,j}(l)|\le C_j(1+|l|^{j+1}),\qquad j\le k. \tag{27}
\]
Induction also gives locally uniform convergence
$F_{\epsilon,j}\to F_j$. For the integral at step $j+1$, truncate
$u$ to a compact interval, use local uniform convergence there, and
bound the tail by a constant times $1+|u|^{j+2}$. Its uniform
integrability follows from the assumed Wasserstein convergence, since
$j+2\le k+1$. The same argument and (27), using a coupling that
converges in $L^{k+1}$, show
$F_{\epsilon,k}(L_\epsilon)\to F_k(L)$ in $L^1$.
Taking half the integrated absolute residual proves the limit in (24).

It remains to show that the limiting constant is positive. Each $F_j$
is continuous and grows at most like $1+|l|^{j+1}$. If
$b=\operatorname{ess\,sup}L<\infty$, then $b>0$, $a(b)=b$, and
the integral term in (23) vanishes at $b$. Hence $F_j(b)=b^{j+1}>0$;
continuity and the definition of essential supremum give positive
probability of $F_j(L)>0$. If the essential supremum is infinite,
induction in (23) gives $F_j(l)/l^{j+1}\to1$ as $l\to+\infty$:
$a(l)/l\to1$, and the integral over $L>l$ tends to zero by the
available moment of order $j+2$ at step $j+1$. Again $F_j(L)>0$ on
an event of positive probability. This proves $C_k>0$.
\end{proof}

\begin{corollary}[Every order before the sector floor is sharp]
For the neural posterior in Theorem~\ref{thm:rate}, every integer
$0\le k\le r-5$ has
\[
             \|\widehat\Pi_{N,k}-\Pi_N\|_{\TV}
                         \sim C_kN^{-(k+1)/2},          \tag{28}
\]
where (23)--(24) use the explicit variable $L=\mathcal L$ in (18).
\end{corollary}
\begin{proof}
The proof of Theorem~\ref{thm:constants} gives weak convergence of
$\sqrt N(w_N-1)$ to $\mathcal L$ and bounded moments of order
$p_*=r-3$. Thus it gives Wasserstein convergence at every smaller
order, including $k+1\le r-4$. Apply Theorem~\ref{thm:hierarchy} with
$\epsilon=N^{-1/2}$. The omitted sector has probability
$O(N^{-p_*/2})=o(N^{-(k+1)/2})$. At $k=r-4$ the missing sector is
already of the same order as the bound, and this perturbation argument
does not identify its leading constant.
\end{proof}

\section{Completion and limitations}
Apply Lemma~\ref{lem:imh} with $p=k+1\le p_*$. Equation (16) gives
$O(N^{-(k+1)/2})$ distance to the canonical conditional posterior.
Uniform random injection is the same Markov operation for the two laws,
so it cannot increase total variation. Add (3) to compare with the full
posterior; $b=p_*/2$ gives (5). At $k=r-4=p_*-1$ this yields
$O(N^{-b})$. Further fixed updates cannot increase distance to the
conditional stationary law. Every output remains in $G_N$, giving the
lower bound $\Pi_N(G_N^c)\ge cN^{-b}$ and hence (6).

At fixed dimension each proposal, output evaluation, and importance
ratio uses $O(N)$ arithmetic operations; there are only $k+1$ proposals
per independent chain. The proposal's positive limiting acceptance
proves the expected-cost statement. No assertion about rapid mixing
from an arbitrary initial network is involved.

The theorem concerns full labels and a width-scaled likelihood. It does
not provide polynomial-in-dimension pretraining, nor does it cover a
hard-margin posterior using the earlier approximate-slack proposal:
that proposal has a support mismatch near the margin boundary. A
different full Gaussian disintegration is needed for the hard case.
Floating-point importance diagnostics are useful checks of the formulas,
but are not certificates of total variation to the full neural posterior.
The proof above is conventional and not fully Lean formalized.

\begin{thebibliography}{9}
\bibitem{Liu} J. S. Liu, \emph{Metropolized independent sampling with
comparisons to rejection sampling and importance sampling}, Statistics
and Computing 6 (1996), 113--119.
\bibitem{DJW} G. Deligiannidis, P. E. Jacob, E. M. Khribch, and G. Wang,
\emph{On importance sampling and independent Metropolis--Hastings with an
unbounded weight function}, arXiv:2411.09514v2 (2025).
\url{https://arxiv.org/abs/2411.09514}.
\end{thebibliography}
\end{document}
